\inputencoding{utf8}
\HeaderA{sap\_day}{Calculate ideal sap days from temperature rasters}{sap.Rul.day}
%
\begin{Description}
Maple sap flows the most on days where temperatures are oscillating between
freezing and thawing, known as the freeze/thaw cycle. This function uses
tmax and tmin rasters (potentially loaded in from \code{\LinkA{loca\_t\_rast()}{loca.Rul.t.Rul.rast}} or
\code{\LinkA{prism\_t\_rast()}{prism.Rul.t.Rul.rast}}) to calculate the proportion and sum of days in each year
that are ideal for maple sap tapping. Using a 30-year record of sap flow
and weather observations at the Uihlein Forest, Lake Placid, NY, sap flow
occurred on 80\% of the days when the minimum temperature fell below
-1.1 °C and maximum temperature exceeded 2.2 °C, which is where the
default values for the temperature threshold come from (Skinner et al.
2010).
\end{Description}
%
\begin{Usage}
\begin{verbatim}
sap_day(tmax_rast, tmin_rast, t_upper = 2.2, t_lower = -1.1, years = NULL)
\end{verbatim}
\end{Usage}
%
\begin{Arguments}
\begin{ldescription}
\item[\code{tmax\_rast}] terra raster stack of tmax values that must have terra::time
values containing date information to subset by year

\item[\code{tmin\_rast}] terra raster stack of tmin values that must have terra::time
values containing date information to subset by year

\item[\code{t\_upper}] upper temperature value for the freeze/thaw cycle that
will be used in a logical statement to find ideal sap tapping days

\item[\code{t\_lower}] lower temperature value for the freeze/thaw cycle that
will be used in a logical statement to find ideal sap tapping days

\item[\code{years}] integer vector of years to subset the raster stacks by. All
of the years must be present in the raster stacks
\end{ldescription}
\end{Arguments}
%
\begin{Value}
list of length two - a raster stack of the proportion of ideal
sap tapping days per year, and a raster stack of the sum of ideal sap
tapping days per year
\end{Value}
%
\begin{Examples}
\begin{ExampleCode}
## Not run: 
# Load in rasters
loca_rast <- loca_t_rast("D:/Data/LOCA2/ACCESS-CM2/0p0625deg/r1i1p1f1/")

# Calculate yearly sap days
k_upper <- 2.2 + 273.15
k_lower <- -1.1 + 273.15
loca_sap_day <- sap_day(loca_rast$tmax, loca_rast$tmin,
                        t_upper = k_upper, t_lower = k_lower)

## End(Not run)
\end{ExampleCode}
\end{Examples}
